\section*{Erd\H{o}s Problem 188: unit-step progressions and the unrestricted-step distinction}

\subsection*{1. Statement and scope}
What is the smallest $k$ for which the plane admits a red-blue colouring with no red pair at distance one and no blue collinear $k$-term progression of step length one? The minimum remains unresolved in this attempt. We prove directly that $k\ge4$ and prove the impossibility of any finite $k$ if the blue step length is unrestricted. We also distinguish our proved lower bound from the stronger established lower bound $k\ge6$.

\subsection*{2. A complete obstruction to avoiding blue triples}
Suppose a colouring has neither a red unit-distance pair nor a blue unit-step triple. If there are no red points, any three points in a unit-step progression contradict the second condition, so choose a red point and translate it to the origin. Every point on the unit circle is blue.

Choose unit vectors $u,v$ with $|u-v|=1$, equivalently $u\cdot v=1/2$. The points $2u-v,u,v$ form a unit-step collinear triple. Its latter two points are blue, so $2u-v$ must be red. Its squared norm is
\[
 |2u-v|^2=4+1-4(u\cdot v)=3.
\]
Rotating the chosen pair about the origin shows that every point on the circle of radius $\sqrt3$ must be red: every point of that circle is a rotation of $2u-v$. But this circle contains the two points
\[
 (-1/2,\sqrt{11}/2),\qquad(1/2,\sqrt{11}/2),
\]
which are at distance one. This contradiction proves that every colouring with no red unit pair has a blue unit-step triple. Thus the sought minimum, if it exists, is at least four. This argument includes no regularity or measurability assumption on the colouring.

The stronger known theorem of Tsaturian states that such a colouring always has a blue unit-step progression of length five, and therefore the known lower bound is $k\ge6$. That theorem is explicit external context; its proof is not supplied by the shorter radius-$\sqrt3$ argument above. The primary paper's statement has been checked, and the older draft's unsupported stronger unpublished claim is not adopted.

\subsection*{3. Why the step-length qualification changes the question}
Assume only that there is no red unit pair and restrict to the integer lattice points on any fixed unit-spaced line. For each integer $n$, at least one of $n,n+1$ is blue. Choose $j(n)\in\{0,1\}$ with $n+j(n)$ blue, choosing zero if both choices work.

For any desired length $r$, van der Waerden's theorem, the explicit standard input here, supplies an arithmetic progression
\[
 a,a+d,\ldots,a+(r-1)d
\]
on which $j(n)$ is constant, say equal to $j$. Translating the progression by $j$ gives a blue progression of length $r$. Its step $d$ is a positive integer, but need not equal one. Since $r$ was arbitrary, forbidding blue progressions of arbitrary step length is impossible for every finite $k$.

This proves the observation attributed to Alon on the problem page. It does not prove that arbitrarily long unit-step blue progressions must occur: that would discard the step restriction on which the actual problem depends.

\subsection*{4. Assessment and sources}
The radius argument and the unrestricted-step conclusion are fully proved. No colouring avoiding both a red unit pair and a sufficiently long unit-step blue progression has been constructed here, and the exact minimum is not determined.

Sources: \url{https://www.erdosproblems.com/188}; S. Tsaturian, \emph{A Euclidean Ramsey result in the plane}, \url{https://arxiv.org/abs/1703.10723}, for the stronger stated five-point theorem. Van der Waerden's theorem is used only in Section 3.

\textbf{Completion rate estimate: 25\%.}

